
\documentclass[11pt,a4paper]{report}

\usepackage[utf8]{inputenc}
\usepackage[T1]{fontenc}
\usepackage[italian]{babel}
\usepackage{textcomp}
\usepackage{siunitx}
\usepackage{caption}
\usepackage{subfig}
\usepackage{graphicx}
\usepackage{amsmath,amssymb}
\usepackage{natbib}
\usepackage{braket}
\usepackage{float}
\usepackage[most]{tcolorbox}
\usepackage{geometry}
\usepackage{hyperref}

\geometry{margin=2.5cm}

\title{QUBO Penalty Calibration}
\author{Maicol Nicolini}
\date{\today}

\begin{document}

\maketitle
\newpage
\tableofcontents
\newpage

% ============================================================
\chapter{Introduzione}
% ============================================================

Consideriamo una formulazione QUBO del tipo

\[
H(\mathbf z)
=
H_{\mathrm{obj}}(\mathbf z)
+
\sum_r \lambda_r H_r(\mathbf z),
\]

dove $H_{\mathrm{obj}}$ rappresenta la funzione obiettivo del problema
originale e i termini $H_r$ codificano i vincoli tramite penalità.

La scelta dei coefficienti $\lambda_r$ è cruciale:
una penalità troppo piccola può rendere conveniente la violazione del
vincolo, mentre una penalità eccessivamente grande può peggiorare la
scalatura numerica del QUBO e, nel caso di solver quantum o
quantum-inspired, alterare sfavorevolmente la scala energetica del
problema.

L'obiettivo di questo documento è raccogliere due famiglie di strategie
per la calibrazione delle penalty:

\begin{enumerate}
    \item \textbf{problem-based penalty calibration}, in cui si sfrutta
    direttamente la struttura del problema e l'interpretazione fisica
    dei termini della funzione obiettivo;
    \item \textbf{relaxation-based penalty calibration}, in cui si
    costruisce un lower bound della funzione obiettivo mediante
    Semidefinite Programming (SDP) relaxation e si combina tale bound con
    una soluzione feasible nota.
\end{enumerate}

Nel seguito utilizzeremo come esempio guida il problema di Wind Farm
Layout Optimization (WFLO) formulato come QUBO.

% ============================================================
\chapter{Problem-based penalty calibration}
% ============================================================

\section{Formulazione QUBO del WFLO}

La formulazione considerata è

\[
H(\mathbf z)
=
H_{\mathrm{wake}}(\mathbf z)
+
\lambda_N
\left(
\sum_i z_i-K
\right)^2
+
\lambda_S
\sum_{(i,j)\in\mathcal E_{\mathrm{inv}}}
z_i z_j,
\]

dove

\[
H_{\mathrm{wake}}(\mathbf z)
=
\sum_{i<j}L_{ij}z_i z_j.
\]

Il coefficiente $L_{ij}\geq 0$ rappresenta la perdita energetica pairwise
dovuta all'interazione aerodinamica tra le turbine nelle candidate
locations $i$ e $j$.

Nel caso del benchmark considerato,

\[
K=81.
\]

Si assume di utilizzare la formulazione QUBO ridotta, nella quale il
termine lineare

\[
-A_0\sum_i z_i
\]

è stato rimosso. Tale termine è infatti costante nello spazio delle
soluzioni feasible, poiché ogni soluzione valida contiene esattamente
$K=81$ turbine.

\section{Principio generale}

L'idea utilizzata per dimensionare le penalty è rispondere, per ciascun
vincolo, alla domanda:

\begin{center}
\emph{Qual è il massimo beneficio che il modello potrebbe ottenere
violando questo vincolo?}
\end{center}

Una volta stimato tale beneficio, si sceglie una penalità sufficientemente
grande da rendere la violazione energeticamente sconveniente.

In generale, si desidera una penalty:

\begin{itemize}
    \item sufficientemente grande da favorire o garantire la feasibility;
    \item sufficientemente piccola da non dominare inutilmente la funzione
    obiettivo.
\end{itemize}

% ------------------------------------------------------------
\section{Vincolo sulla cardinalità}
% ------------------------------------------------------------

Il problema impone

\[
\sum_i z_i = K.
\]

Poiché $L_{ij}\geq 0$, in assenza di questo vincolo il modello sarebbe
naturalmente incentivato a selezionare meno turbine, riducendo il numero
di interazioni aerodinamiche.

Il problema consiste quindi nello stimare quanto possa diminuire

\[
H_{\mathrm{wake}}(\mathbf z)
\]

rimuovendo una turbina.

\subsection{Stima conservativa sulla candidate grid}

Consideriamo una generica candidate location $i$.

Se la turbina $i$ viene rimossa, scompaiono tutte le interazioni che la
coinvolgono:

\[
L_{ij}z_i z_j.
\]

In un layout composto da $K$ turbine, la turbina $i$ può interagire con
al massimo $K-1$ turbine.

Definiamo quindi

\[
\mathcal J_i^{K-1}
=
\operatorname{Top}_{K-1}
\left\{
L_{ij}:j\neq i
\right\}.
\]

Nel nostro caso,

\[
K-1=80.
\]

Una stima conservativa del massimo beneficio ottenibile rimuovendo la
turbina $i$ è

\[
m_i^{\mathrm{grid}}
=
\sum_{j\in\mathcal J_i^{K-1}}
L_{ij}.
\]

Questa quantità assume che la turbina $i$ sia contemporaneamente
selezionata con le $K-1$ candidate locations che producono le maggiori
wake losses.

Iterando il procedimento per ogni $i$ si ottiene

\[
m_{\max}^{\mathrm{grid}}
=
\max_i
m_i^{\mathrm{grid}}.
\]

Il termine di cardinalità è

\[
H_N(\mathbf z)
=
\lambda_N
\left(
\sum_i z_i-K
\right)^2.
\]

Rimuovendo una sola turbina:

\[
(K-1-K)^2=1,
\]

quindi la penalità pagata è esattamente

\[
\lambda_N.
\]

Per rendere la rimozione non conveniente rispetto al massimo beneficio
stimato, scegliamo

\[
\lambda_N
>
m_{\max}^{\mathrm{grid}}.
\]

Introducendo un safety factor $\alpha_N>1$:

\[
\boxed{
\lambda_N
=
\alpha_N
m_{\max}^{\mathrm{grid}}
}
\]

con, ad esempio,

\[
\alpha_N\in[1.2,1.8].
\]

Questa stima è prudenziale poiché considera un worst-case costruito
sull'intero spazio delle candidate locations.

\subsection{Calibrazione mediante una soluzione feasible di riferimento}

Se è disponibile un layout feasible di riferimento composto esattamente
da $K$ turbine, si può ottenere una stima meno conservativa.

Sia

\[
\mathbf z^{\mathrm{ref}}
\in
\{0,1\}^{N}
\]

una soluzione feasible tale che

\[
\sum_i z_i^{\mathrm{ref}}=K.
\]

Definiamo

\[
\mathcal S_{\mathrm{ref}}
=
\left\{
i:z_i^{\mathrm{ref}}=1
\right\},
\]

con

\[
|\mathcal S_{\mathrm{ref}}|=K.
\]

Per ciascuna turbina $i\in\mathcal S_{\mathrm{ref}}$, la wake loss
marginale è

\[
m_i^{\mathrm{ref}}
=
\sum_{\substack{
j\in\mathcal S_{\mathrm{ref}}\\
j\neq i
}}
L_{ij}.
\]

Equivalentemente,

\[
m_i^{\mathrm{ref}}
=
\sum_{j\neq i}
L_{ij}z_j^{\mathrm{ref}}.
\]

La quantità $m_i^{\mathrm{ref}}$ rappresenta esattamente la riduzione
della pairwise wake-loss objective che si otterrebbe rimuovendo la turbina
$i$ dal layout di riferimento.

Definiamo

\[
m_{\max}^{\mathrm{ref}}
=
\max_{i\in\mathcal S_{\mathrm{ref}}}
m_i^{\mathrm{ref}}.
\]

La cardinality penalty viene quindi calibrata come

\[
\boxed{
\lambda_N
=
\alpha_N
m_{\max}^{\mathrm{ref}}
}
\]

con $\alpha_N>1$.

In generale ci aspettiamo

\[
m_{\max}^{\mathrm{ref}}
\leq
m_{\max}^{\mathrm{grid}},
\]

poiché la soluzione di riferimento contiene interazioni effettivamente
presenti in un layout feasible, mentre la stima grid-based combina
interazioni worst-case che potrebbero non essere simultaneamente
realizzabili.

La soluzione feasible può essere, ad esempio:

\begin{itemize}
    \item il baseline fornito dal benchmark;
    \item una soluzione ottenuta con OR-Tools;
    \item una soluzione prodotta da una greedy heuristic;
    \item un'altra soluzione feasible nota.
\end{itemize}

Questa seconda procedura è una calibrazione empirica e locale: fornisce
una stima più rappresentativa di un layout realistico, ma non costituisce
necessariamente un bound globale valido per ogni configurazione della
candidate grid.

% ------------------------------------------------------------
\section{Vincolo sulla distanza minima}
% ------------------------------------------------------------

Definiamo l'insieme delle coppie non ammissibili come

\[
\mathcal E_{\mathrm{inv}}
=
\left\{
(i,j):
\|p_i-p_j\|<d_{\min}
\right\}.
\]

Nel benchmark considerato,

\[
d_{\min}=396\text{ m}.
\]

Il termine di penalità è

\[
H_S(\mathbf z)
=
\lambda_S
\sum_{(i,j)\in\mathcal E_{\mathrm{inv}}}
z_i z_j.
\]

Se una coppia non ammissibile viene selezionata, la funzione obiettivo
riceve una penalità pari a $\lambda_S$.

La differenza rispetto alla cardinalità è che non è semplice associare il
beneficio della violazione alla sola wake loss $L_{ij}$ della coppia
invalida.

Consentire una coppia non feasible può infatti modificare l'intera
configurazione del layout e permettere al solver di ottenere una
wake-loss objective complessivamente inferiore.

\subsection{Stima conservativa sulla candidate grid}

Un layout di $K$ turbine contiene esattamente

\[
\binom{K}{2}
\]

interazioni pairwise.

Con $K=81$:

\[
\binom{81}{2}=3240.
\]

Definiamo

\[
\mathcal L^{\binom{K}{2}}
=
\operatorname{Top}_{\binom{K}{2}}
\left\{
L_{ij}:i<j
\right\}.
\]

Nel caso specifico:

\[
\mathcal L^{3240}
=
\operatorname{Top}_{3240}
\left\{
L_{ij}:i<j
\right\}.
\]

Definiamo quindi

\[
E_{\max}^{\mathrm{grid}}
=
\sum_{L_{ij}\in\mathcal L^{3240}}
L_{ij}.
\]

Questa quantità rappresenta un upper bound conservativo sulla wake-loss
objective di qualunque configurazione di 81 turbine.

È un bound conservativo perché le 3240 interazioni selezionate
globalmente potrebbero coinvolgere più di 81 candidate locations e quindi
non essere simultaneamente realizzabili all'interno di una singola
configurazione.

Per rendere una singola spacing violation più costosa di tale scala
energetica si può scegliere

\[
\lambda_S
>
E_{\max}^{\mathrm{grid}}.
\]

Con safety factor $\alpha_S>1$:

\[
\boxed{
\lambda_S
=
\alpha_S
E_{\max}^{\mathrm{grid}}
}
\]

con, ad esempio,

\[
\alpha_S\in[1.1,1.5].
\]

Questa strategia è robusta ma spesso estremamente conservativa.

\subsection{Calibrazione mediante una soluzione feasible di riferimento}

Sia $\mathbf z^{\mathrm{ref}}$ una soluzione feasible contenente
esattamente $K$ turbine e rispettante il vincolo di distanza minima.

Definiamo la sua pairwise wake-loss objective:

\[
E_{\mathrm{ref}}
=
H_{\mathrm{wake}}(\mathbf z^{\mathrm{ref}})
=
\sum_{i<j}
L_{ij}
z_i^{\mathrm{ref}}
z_j^{\mathrm{ref}}.
\]

Equivalentemente,

\[
E_{\mathrm{ref}}
=
\sum_{\substack{
i<j\\
i,j\in\mathcal S_{\mathrm{ref}}
}}
L_{ij}.
\]

Si sceglie la spacing penalty in modo che una singola violazione abbia
un costo superiore all'intera wake-loss objective della soluzione
feasible nota:

\[
\lambda_S
>
E_{\mathrm{ref}}.
\]

Introducendo un safety factor $\alpha_S>1$:

\[
\boxed{
\lambda_S
=
\alpha_S
E_{\mathrm{ref}}
}
\]

con, ad esempio,

\[
\alpha_S=1.2.
\]

Infatti, per una soluzione infeasible $\mathbf z^{\mathrm{inf}}$ con
almeno una spacing violation,

\[
H_S(\mathbf z^{\mathrm{inf}})
\geq
\lambda_S.
\]

Poiché

\[
H_{\mathrm{wake}}(\mathbf z^{\mathrm{inf}})
\geq0,
\]

segue che

\[
H(\mathbf z^{\mathrm{inf}})
\geq
\lambda_S.
\]

Se

\[
\lambda_S
>
E_{\mathrm{ref}},
\]

allora

\[
H(\mathbf z^{\mathrm{inf}})
>
E_{\mathrm{ref}}
=
H_{\mathrm{wake}}(\mathbf z^{\mathrm{ref}}).
\]

Quindi una configurazione contenente almeno una spacing violation non può
battere la soluzione feasible di riferimento.

Anche in questo caso il valore ottenuto può essere sottoposto a
sensitivity analysis, riducendo progressivamente $\lambda_S$ e
verificando empiricamente la feasibility delle soluzioni prodotte.

% ============================================================
\chapter{Relaxation-based penalty calibration}
% ============================================================

\section{Il problema Big-M}

Consideriamo un problema binario quadratico vincolato:

\[
\min_{\mathbf x\in\{0,1\}^n}
f(\mathbf x)
=
\mathbf x^TQ\mathbf x
\]
soggetto a
\[
A\mathbf x=\mathbf b.
\]

Per trasformarlo in un QUBO, il vincolo viene promosso a termine di
penalità:

\[
H_M(\mathbf x)
=
f(\mathbf x)
+
M\|A\mathbf x-\mathbf b\|^2.
\]

Per una soluzione feasible la penalty è nulla. Per una soluzione
infeasible è positiva. \newline

Il valore $M$ deve quindi essere sufficientemente grande affinché nessuna
soluzione infeasible possa avere energia inferiore alla soluzione ottima
feasible. \newline

Una scelta eccessivamente grande di $M$ introduce però il cosiddetto
\emph{Big-M problem}: la funzione obiettivo viene dominata dalla penalty,
peggiorando la scala numerica dei coefficienti e, nel contesto quantum e
quantum-inspired, la scala energetica dell'Hamiltoniana.
L'obiettivo è quindi scegliere la più piccola penalty che renda la
riformulazione esatta.

\section{Upper bound da una soluzione feasible}
Sia
\[
\mathbf x_{\mathrm{feas}}
\]
una soluzione feasible nota. Il suo valore
\[
U
=
f(\mathbf x_{\mathrm{feas}})
\]
costituisce un upper bound sul valore ottimo del problema vincolato:
\[
f(\mathbf x^\star)
\leq
U.
\]

La soluzione feasible può provenire da:

\begin{itemize}
    \item un solver classico con time limit limitato;
    \item una greedy heuristic;
    \item una soluzione benchmark;
    \item una soluzione precedentemente calcolata.
\end{itemize}

Più la soluzione feasible è buona, più piccolo sarà $U$ e più stretta
sarà la penalty ottenuta.

\section{Lower bound della funzione obiettivo}
Rimuoviamo temporaneamente i vincoli e consideriamo
\[
\min_{\mathbf x\in\{0,1\}^n}
f(\mathbf x).
\]
Supponiamo di conoscere un lower bound $L$ tale che

\[
L
\leq
\min_{\mathbf x\in\{0,1\}^n}
f(\mathbf x).
\]
La differenza
\[
U-L
\]
rappresenta una stima conservativa del massimo vantaggio che una soluzione
potrebbe ottenere ignorando i vincoli. \newline

Per rendere sconveniente tale vantaggio è sufficiente scegliere
\[
M>U-L.
\]

Introducendo un margine $\delta>0$:
\[
\boxed{
M
=
U-L+\delta
}
\]
Il lower bound viene trovata tramite una procedura di \emph{relaxation}.
Ne esistono diverse, le più importanti sono due: \emph{linear relaxation} e 
\emph{Semidefinite Programming relaxation}, ci concentriamo su quest'ultima. 
\[
\boxed{
M_{\mathrm{SDP}}
=
f(\mathbf x_{\mathrm{feas}})
-
f^\star_{\mathrm{SDP}}
+
\delta
}
\]

\section{Teoria della SDP relaxation}

Consideriamo una funzione binaria quadratica generale:
\[
f(\mathbf x)
=
\mathbf x^TQ\mathbf x
+
\mathbf L^T\mathbf x,
\qquad
\mathbf x\in\{0,1\}^n.
\]
L'obiettivo è ottenere un lower bound senza risolvere esattamente il
problema combinatorio.

\subsection{Lifting}

Definiamo il vettore aumentato
\[
\mathbf y
=
\begin{pmatrix}
1\\
\mathbf x
\end{pmatrix}
\]
e la matrice lifted
\[
Y
=
\mathbf y\mathbf y^T
=
\begin{pmatrix}
1 & \mathbf x^T\\
\mathbf x & \mathbf x\mathbf x^T
\end{pmatrix}.
\]

La parte in basso a destra contiene direttamente i prodotti pairwise:

\[
Y_{i+1,j+1}
=
x_i x_j.
\]

Definendo
\[
\widetilde Q
=
\begin{pmatrix}
0 & \frac{1}{2}\mathbf L^T\\
\frac{1}{2}\mathbf L & Q
\end{pmatrix},
\]
la funzione obiettivo può essere riscritta come
\[
f(\mathbf x)
=
\operatorname{Tr}
\left(
Y^T\widetilde Q
\right).
\]
Dopo il lifting, la funzione obiettivo è quindi lineare nella matrice
$Y$.

\subsection{Formulazione esatta}

Se
\[
Y
=
\mathbf y\mathbf y^T,
\]

allora $Y$ soddisfa due proprietà:
\[
Y\succeq0
\]

e
\[
\operatorname{rank}(Y)=1.
\]
La condizione
\[
Y\succeq0
\]
significa che $Y$ è positiva semidefinita, cioè
\[
\mathbf v^T Y \mathbf v
\geq0
\]

per qualsiasi vettore $\mathbf v$.  

Infatti,
\[
\mathbf v^TY\mathbf v
=
\mathbf v^T\mathbf y\mathbf y^T\mathbf v
=
(\mathbf y^T\mathbf v)^2
\geq0.
\]

Il vincolo problematico è
\[
\operatorname{rank}(Y)=1,
\]

che rende il problema non convesso.

\subsection{Relaxation}

La SDP relaxation rimuove il vincolo di rango 1 e mantiene la condizione
semidefinita positiva. Inoltre, poiché per una variabile binaria vale
\[
x_i^2=x_i,
\]
si impone la relazione lineare
\[
Y_{1,i+1}
=
Y_{i+1,i+1}.
\]
Si aggiungono inoltre bounds del tipo
\[
0\leq Y_{ij}\leq1.
\]
La relaxation assume quindi la forma
\[
\min_Y
\operatorname{Tr}
\left(
Y^T\widetilde Q
\right)
\]
soggetta a
\[
Y\succeq0,
\]

\[
Y_{1,i+1}
=
Y_{i+1,i+1},
\]

\[
0\leq Y_{ij}\leq1.
\]

\section{Perché la SDP fornisce un lower bound}
Nel problema originale la matrice $Y$ deve necessariamente essere
generata da un vettore binario e avere rango 1. Nella SDP ammettiamo anche matrici positive semidefinite che non
corrispondono ad alcuna soluzione binaria reale. \newline
Abbiamo quindi
\[
\mathcal F_{\mathrm{binary}}
\subseteq
\mathcal F_{\mathrm{SDP}}.
\]
Poiché stiamo minimizzando,
\[
\boxed{
f^\star_{\mathrm{SDP}}
\leq
f^\star_{\mathrm{binary}}
}
\]
e il valore ottimo della SDP è quindi un lower bound valido per il
problema binario originale.

\section{Interpretazione della penalty SDP}
Una volta disponibili:
\begin{itemize}
    \item una soluzione feasible con costo
    \[
    U=f(\mathbf x_{\mathrm{feas}});
    \]
    \item un lower bound
    \[
    f^\star_{\mathrm{SDP}};
    \]
\end{itemize}

si definisce
\[
\boxed{
M_{\mathrm{SDP}}
=
U
-
f^\star_{\mathrm{SDP}}
+
\delta.
}
\]

L'interpretazione è:
\[
\underbrace{
f^\star_{\mathrm{SDP}}
}_{\text{scenario ottimistico}}
\qquad
\longrightarrow
\qquad
\underbrace{
f(\mathbf x_{\mathrm{feas}})
}_{\text{soluzione feasible nota}}.
\]
La differenza tra questi due valori rappresenta una stima conservativa
del massimo vantaggio che una soluzione potrebbe ottenere ignorando i
vincoli. \newline
La SDP viene quindi utilizzata come preprocessing classico per ottenere
una penalty più stretta, non come solver finale del QUBO.

\section{Workflow relaxation-based}

Il procedimento può essere riassunto come segue:

\begin{enumerate}
    \item trovare una soluzione feasible $\mathbf x_{\mathrm{feas}}$;
    \item calcolare
    \[
    U=f(\mathbf x_{\mathrm{feas}});
    \]
    \item rimuovere temporaneamente i vincoli;
    \item risolvere la SDP relaxation;
    \item ottenere
    \[
    f^\star_{\mathrm{SDP}};
    \]
    \item scegliere
    \[
    M_{\mathrm{SDP}}
    =
    U-f^\star_{\mathrm{SDP}}+\delta;
    \]
    \item costruire il QUBO finale con la penalty ottenuta;
    \item validare la feasibility con il solver target.
\end{enumerate}

% ============================================================
\chapter{Confronto tra i due approcci}
% ============================================================

Le due strategie rispondono allo stesso obiettivo: ottenere penalty
sufficientemente grandi da rendere sconveniente la violazione dei
vincoli, evitando però valori inutilmente elevati.

\section{Problem-based}

L'approccio problem-based sfrutta direttamente:

\begin{itemize}
    \item la struttura della funzione obiettivo;
    \item l'interpretazione fisica dei termini;
    \item informazioni specifiche sul vincolo;
    \item eventuali soluzioni feasible note.
\end{itemize}

Vantaggi:

\begin{itemize}
    \item semplice da interpretare;
    \item spesso economico computazionalmente;
    \item può sfruttare fortemente la struttura specifica del problema.
\end{itemize}

Limiti:

\begin{itemize}
    \item le formule dipendono dal problema;
    \item i bound conservativi possono essere molto grandi;
    \item le versioni basate su una reference solution sono in genere
    empiriche e locali.
\end{itemize}

\section{Relaxation-based}

L'approccio SDP costruisce invece una ricetta più generale:

\[
M_{\mathrm{SDP}}
=
U-L+\delta.
\]

Vantaggi:

\begin{itemize}
    \item formulazione generica;
    \item lower bound matematicamente motivato;
    \item spesso produce valori Big-M più stretti rispetto a bound
    elementari.
\end{itemize}

Limiti:

\begin{itemize}
    \item richiede la soluzione di una SDP;
    \item la SDP densa può diventare costosa per problemi di grandi
    dimensioni;
    \item la qualità della penalty dipende dalla qualità sia
    dell'upper bound feasible sia del lower bound SDP.
\end{itemize}

\section{Strategia consigliata}

Una strategia pratica può combinare i due approcci:
\[
\boxed{
\text{problem-based bound}
\rightarrow
\text{SDP tightening}
\rightarrow
\text{sensitivity analysis}
\rightarrow
\text{solver target}
}
\]
Il bound problem-based fornisce una prima scala interpretabile.
La SDP relaxation consente poi di restringere la penalty in modo più
sistematico.
Infine, una sensitivity analysis permette di verificare empiricamente
che la penalty scelta rimanga sufficientemente robusta sul solver target.

% ============================================================
\chapter{Riferimenti}
% ============================================================

E. Alessandroni, S. Ramos-Calderer, I. Roth, E. Traversi e L. Aolita,
\emph{Alleviating the quantum Big-M problem},
arXiv:2307.10379v4, 2025.

\end{document}
